%	RESUMEN
\chapter{RESUMEN}

El crecimiento sostenido del tr\'afico digital y la sofisticaci\'on de los ataques cibern\'eticos han convertido a los centros de datos en objetivos prioritarios para los actores maliciosos, siendo el ransomware, la denegaci\'on de servicio distribuida y el movimiento lateral las t\'ecnicas de mayor impacto sobre la continuidad operativa de las instituciones de educaci\'on superior. El presente proyecto integrador propone el desarrollo de una plataforma de ciberseguridad que integra la detecci\'on de intrusiones basada en aprendizaje autom\'atico con la capacidad de respuesta automatizada sobre redes programables, articulando la detecci\'on de amenazas con la mitigaci\'on de las mismas en la misma plataforma. El sistema se desarroll\'o sobre una arquitectura de cuatro capas: percepci\'on de tr\'afico en el borde de la red, an\'alisis y clasificaci\'on de eventos mediante un modelo de aprendizaje autom\'atico basado en \'arboles de gradiente boosting, control de la infraestructura mediante un controlador de red programable y una capa de visualizaci\'on web que centraliza la monitorizaci\'on y la gesti\'on de la informaci\'on. Se implement\'o un entorno virtual de pruebas con tres hosts y un switch programable para la validaci\'on controlada de los flujos de ataque. Los resultados experimentales demostraron que la plataforma identifica y visualiza los eventos de seguridad en tiempo real y permite la mitigaci\'on de las conexiones maliciosas desde una interfaz \'unica de administraci\'on, proporcionando una base s\'olida para su despliegue en el centro de datos de la FIEC.

\textbf{Palabras Clave:} Ciberseguridad, Redes Programables, Detecci\'on de Intrusiones, Inteligencia Artificial, Centro de Datos


%	ABSTRACT
\chapter{ABSTRACT}

The sustained growth of digital traffic and the increasing sophistication of cyberattacks have turned data centers into prime targets for malicious actors, with ransomware, distributed denial of service and lateral movement being the techniques that most severely affect the operational continuity of higher education institutions. This project develops a cybersecurity platform that integrates machine-learning-based intrusion detection with automated response over software-defined networks, articulating threat detection and mitigation within the same platform. The system was implemented on a four-layer architecture: traffic perception at the network edge, event analysis and classification through a gradient boosting decision tree model, infrastructure control through a programmable network controller, and a web visualization layer that centralizes monitoring and information management. A virtual test environment with three hosts and a programmable switch was used for controlled validation of attack flows. Experimental results showed that the platform identifies and visualizes security events in real time and allows malicious connections to be mitigated from a single administration interface, providing a solid foundation for its deployment at the FIEC data center.

\textbf{Keywords:} Cybersecurity, Software-Defined Networking, Intrusion Detection, Artificial Intelligence, Data Center
